\subsection{Study Limitations and Future Work}

While the proposed CV system establishment procedure demonstrates promising results in both concentrated and sparse ant images, several noteworthy limitations emerged when testing on external datasets. To illustrate these limitations, we evaluated the system using images from the publicly available iNaturalist dataset \cite{horn_inaturalist_2018}. We selected six most representative images from the dataset to demonstrate the limitation. Of the six test images, four contained Argentine ants (Linepithema humile) (Figure \ref{fig:inat} a, b, c, e) and two featured red harvester ants (Pogonomyrmex barbatus) (Figure \ref{fig:inat} d, f). Since these species differ from those used in our training dataset, they provide a valuable test case for assessing model limitations.
The RT-DETR-L model, identified as the best-performing model in Study 1, was trained on over 1,000 ant images from our experimental dataset. Consequently, the model struggled to accurately detect ants in these unfamiliar images, particularly when subjects were larger than expected or blended with complex backgrounds. The model performed well on relatively clear images (Figure \ref{fig:inat}a and d), but struggled with images containing small debris such as sand or leaf shadows, producing numerous false positives (Figure \ref{fig:inat}b and e). A particularly illustrative example appears in Figure \ref{fig:inat}c, where the ants are significantly larger than those in our calibration set, with much more prominent morphological features. Here, the model failed to detect the actual ants but incorrectly identified ant legs as separate ants due to their similar size and coloration. Finally, in Figure \ref{fig:inat}f, where ants are barely visible even to human observers due to background camouflage, the model again failed to achieve detection.
These examples underscore the model's limitations when applied to images that differ substantially from the calibration dataset. While not the primary focus of this study, broader application would require consideration of ant species diversity and morphological variation. Enhanced model robustness and generalizability could be achieved by incorporating a wider range of ant images in the calibration set, preferably spanning multiple species, complex backgrounds, and varied distance scales.
Several approaches could address these limitations. First, a more complex model architecture, such as RT-DETR-X with 34 million additional parameters (nearly 80\% more than RT-DETR-L), might improve generalizability. However, this approach would require higher-grade GPUs with sufficient video memory for model calibration, more extensive training datasets to prevent overfitting, and longer processing times compared to smaller models.
Alternatively, rather than increasing model complexity or dataset size, a data augmentation approach could simulate debris and background complexity within the existing calibration set. For instance, treating images as two-dimensional matrices, we could initialize multiple multivariate Gaussian distributions to model the probability of debris and background complexity, then randomly sample from these distributions to generate dark or light pixel clusters that mimic natural debris. While this approach may seem simplistic, it could provide a valuable starting point for improving model performance from a different methodological perspective.

\begin{figure}[H]
    \centering
    \includegraphics[width=\textwidth]{figure_12.jpg}
    \caption{Ant images from the iNaturalist dataset \cite{horn_inaturalist_2018}. The ant species are \emph{Linepithema humile} (a, b, c, e) and \emph{Pogonomyrmex barbatus} (d, f). The blue bounding boxes indicate the detection results from the RT-DETR-L model calibrated in Study 1 and make detection by Slicing Aided Hyper Inference (SAHI) \cite{akyon_slicing_2022}.}
    \label{fig:inat}
\end{figure}

\textit{Alt text}: 


\subsection{Data Dissemination and Web-Based Application}

The source code, the dataset organized in YOLO format, along with the calibrated YOLO11n, YOLO11m, and RT-DETR-L model weights (.pt), can be accessed for generating and evaluating the CV system and is publicly available on GitHub (\url{https://github.com/Niche-Lab/ant-detective}). To support researchers and practitioners, a web-based Streamlit application called Ant Detective (\url{https://ant-detective.streamlit.app}) has also been developed. This application allows users to upload images and receive a folder containing detection results in YOLO format and visualized images with blue bounding boxes.
